
正如章题或已表明的，本章包含的是关于本书第二部分的练习。

\section{第~5~章的练习}
\label{sec:ch14-exercises-ch5}

\begin{exercise}\label{ex:ch14-001-symplectic-orthogonal}
设 $E$ 是一个维数为 $2n$ 的\trSymplectic\trVectorSpace。若 $F$ 是 $E$ 的一个子集，我们就以 $F^\circ$ 表示对于\trSymplecticForm{}而言\trOrthogonal{}的向量子空间，即由一切使得 $\omega(x,y)=0$（对每一个 $y\in F$）的向量 $x$ 所组成。

对于 $E$ 的一个子空间 $F$，证明
\[
\dim F+\dim F^\circ=\dim E\text{。}
\]

若\trForm{} $\omega$ 在 $F$ 上为零，也就是说，若 $F\subset F^\circ$，我们就称 $F$ 是\trIsotropic{}的。关于 $F$ 的维数我们能说些什么？若 $F$ 是\trIsotropic{}的且 $\dim F=n$，我们就称 $F$ 是\trLagrange{}式的。

设 $F$ 是一个“\trCoisotropic”子空间（也就是说，其\trOrthogonal{}空间是\trIsotropic{}的）。证明：$\omega$ 在商\trVectorSpace{} $F/F^\circ$ 上诱导出一个\trSymplecticForm。设 $L$ 是 $E$ 的一个\trLagrange{}式子空间，满足 $L+F=E$。证明：复合
\[
L\cap F\subset F\longrightarrow F/F^\circ
\]
是一个\trLagrange{}式子空间的\trInjective。
\end{exercise}

\begin{exercise}\label{ex:ch14-002-liouville-pullback}
设 $V$ 是一个\trManifold，并设 $\eta$ 是 $V$ 上的一个 1-\trForm。我们可以把 $\eta$ 看成\trCotangent{}\trVectorBundle{}的一个\trSection
\[
\eta\colon V\longrightarrow T^\star V\text{，}
\]
若 $\lambda$ 表示 $T^\star V$ 上的\trLiouville{}形式，那么 $V$ 上的\trForm{} $\eta^\star\lambda$ 是什么？
\end{exercise}

\begin{exercise}[\trLagrange{}\trSubmanifold]\label{ex:ch14-003-lagrangian-submanifolds}\BookIndexDefinition{T005}
若 $(W,\omega)$ 是一个维数为 $2n$ 的\trSymplectic\trManifold，那么对于\trSubmanifold{} $j\colon L\subset W$，只要 $\dim L=n$ 且 $j^\star\omega=0$，我们就称它为\emph{\trLagrange{}式的}。

\begin{enumerate}[label=(\arabic*)]
\item 证明：在一个\trSymplectic{}曲面中，所有的曲线都是\trLagrange{}式的。
\item 证明：若 $L_1$ 在 $W_1$ 中是\trLagrange{}式的，而 $L_2$ 在 $W_2$ 中是\trLagrange{}式的，那么 $L_1\times L_2$ 在 $W_1\times W_2$ 中是\trLagrange{}式的。构造 $\RR^{2n}$ 中的一个\trLagrange{}式\trTorus{} $T^n$。证明：每一个\trSymplectic\trManifold{}都含有\trLagrange{}式\trTorus。
\item 设 $\eta$ 是\trManifold{} $V$ 上的一个 1-\trForm。我们把 $\eta$ 看成 $T^\star V$ 的一个\trSection（如\cref{ex:ch14-002-liouville-pullback}中那般）。证明：$\eta$ 的\trImage{}是一个\trLagrange{}式\trSubmanifold{}当且仅当\trForm{} $\eta$ 是\trClosed{}的。
\item 设 $f$ 是\trCompact{}\trManifold{} $V$ 上的一个 $\cC^\infty$ 函数。由上所述，1-\trForm{} $df$ 定义了 $T^\star V$ 的一个\trLagrange{}式\trSubmanifold。证明：它必与\trZeroSection{}相交。
\item 考虑乘积\trManifold{} $W\times W$，并给它配备\trSymplecticForm{} $\omega\oplus(-\omega)$。证明：对角线是一个\trLagrange{}式\trSubmanifold。设 $\varphi$ 是从 $W$ 到其自身的一个\trDiffeomorphism。在 $\varphi$ 满足什么条件时，$\varphi$ 的图像是 $W\times W$ 的一个\trLagrange{}式\trSubmanifold？
\end{enumerate}
\end{exercise}

\begin{exercise}\label{ex:ch14-004-cotangent-integral}
设 $w$ 是从\trSphere{} $S^2$ 到\trManifold{} $V$ 的\trCotangent{}空间 $T^\star V$ 的一个 $\cC^\infty$ 映射。证明\footnote{译注：原书此处径用 $\omega$ 而未重申记号；按第~5~章的口径，$\omega$ 即 $T^\star V$ 上的\trCanonical{}\trSymplecticForm{} $\omega=d\lambda$。——译者注}
\[
\int_{S^2}w^\star\omega=0\text{。}
\]
\end{exercise}

\begin{exercise}\label{ex:ch14-005-symplectic-surface}
设 $\omega$ 是\trCompact{}曲面 $\Sigma$ 上的一个\trSymplecticForm。证明 $\Sigma$ 是\trOrientation{}的，进而证明
\[
\int_{\Sigma}\omega\neq 0\text{。}
\]

考虑由包含关系 $\CC^2\subset\CC^{n+1}$ 所诱导的包含映射 $\CP^1(\CC)\subset\CP^n(\CC)$。证明：$\CP^1(\CC)$ 是一个\trSymplectic\trSubmanifold。由此推出：存在一个 $\cC^\infty$ 映射 $w\colon S^2\longrightarrow\CP^n(\CC)$，使得当 $\omega$ 现在表示 $\CP^n(\CC)$ 的\trSymplecticForm{}时，
\[
\int_{S^2}w^\star\omega\neq 0\text{。}
\]
\end{exercise}

\begin{exercise}\label{ex:ch14-006-graph-diagonal-transverse}
设 $\varphi^1$ 是 $W$ 的一个\trHamiltonian\trDiffeomorphism。证明：$\varphi^1$ 的图像在 $(x,x)$ 处与 $W\times W$ 的对角线\trTransverse，当且仅当 $x$ 的\trTrajectory{}是\trNondegenerate{}的。
\end{exercise}

\begin{exercise}[\trHamiltonian\trDiffeomorphism{}群]\label{ex:ch14-007-hamiltonian-diffeomorphism-group}\BookIndexDefinition{T006}
在本练习中，我们考虑一个配备了\trSymplecticForm{} $\omega$ 的\trCompact\trManifold{} $W$。我们以 $\mathrm{Ham}(W)$ 表示 $W$ 的那样一些\trDiffeomorphism{}所成的集合：它们是\trHamiltonian{}流在时刻 1 处的映射。

\begin{enumerate}[label=(\arabic*)]
\item 假设\trHamiltonian{}量 $H_t$ 生成\trIsotopy{} $\varphi_t$，而\trHamiltonian{}量 $K_t$ 生成 $\psi_t$。令
\[
G_t=H_t+K_t\circ\varphi_t^{-1}\text{。}
\]
证明
\[
X_{G_t}(x)=X_{H_t}(x)+\bigl(T_{\varphi_t^{-1}(x)}\varphi_t\bigr)\bigl(X_{K_t}\circ\varphi_t^{-1}(x)\bigr)\text{。}
\]
由此推出：$\varphi_t\circ\psi_t$ 是由\trHamiltonian\footnote{即使 $H$ 与 $K$ 都是\trAutonomous{}的，复合而成的\trHamiltonian\trIsotopy{}也（一般说来）不来自一个\trAutonomous{}的\trHamiltonian{}量。附加题：何时会如此？}$H_t+K_t\circ\varphi_t^{-1}$ 生成的\trHamiltonian\trIsotopy。
\item 描述由 $-H_t\circ\varphi_t$ 所生成的\trHamiltonian\trIsotopy。
\item 证明：集合 $\mathrm{Ham}(W)$ 是 $W$ 的\trSymplectic{}\trDiffeomorphism{}群的一个子群。
\item 设 $\varPhi$ 是 $W$ 的一个\trSymplectic\trDiffeomorphism。描述由\trHamiltonian{}量 $H_t\circ\varPhi$ 所生成的\trHamiltonian\trIsotopy。证明：子群 $\mathrm{Ham}(W)$ 是 $W$ 的全体\trSymplectic{}\trDiffeomorphism{}群的一个正规子群。
\end{enumerate}
\end{exercise}

\begin{exercise}\label{ex:ch14-008-jacobian-flow}
设 $x$ 是 $\RR^{2n}$ 上一个\trAutonomous{}的\trHamiltonian{}量 $H$ 的一个\trNondegenerate\trCriticalPoint。我们以 $\varphi^t$ 表示 $X_H$ 的\trFlow。试在坐标 $(p_i,q_i)$ 下写出 $X_H$。证明：$\varphi^t$ 的\trJacobi{}矩阵满足
\[
\frac{d}{dt}\Bigl(\operatorname{Jac}_x\varphi^t\Bigr)=\bigl(J_0\operatorname{Hess}_x(H)\bigr)\cdot\operatorname{Jac}_x\varphi^t
\]
从而有
\[
\operatorname{Jac}_x\varphi^t=e^{tJ_0\operatorname{Hess}_x(H)}\text{。}
\]

设 $x$ 是\trSymplectic\trManifold{} $W$ 上一个\trAutonomous{}的\trHamiltonian{}量 $H$ 的一个\trCriticalPoint。证明：若 $x$ 作为\trPeriodic\trTrajectory{}是\trNondegenerate{}的，那么它作为 $H$ 的\trCriticalPoint{}也是\trNondegenerate{}的。
\end{exercise}

\begin{exercise}[谐振子]\label{ex:ch14-009-harmonic-oscillator}
在\trSymplectic\trManifold{} $\RR^{2n}$ 上，我们考虑\trHamiltonian{}量
\[
H(q,p)=\frac{1}{2}\sum_i\alpha_i\bigl(p_i^2+q_i^2\bigr)\text{，}
\]
其中诸 $\alpha_i$ 都是正实数，比如说 $0<\alpha_1\le\alpha_2\le\dotsb\le\alpha_n$。

\begin{enumerate}[label=(\arabic*)]
\item 写出相应的\trHamiltonian{}系统并把它解出来。
\item 我们假设所有的 $\alpha_i/\alpha_j$ 都是无理数。证明：该方程组恰有 $n$ 族\trPeriodic{}解，每一族都含在一个平面之内。我们取定含在方程为 $(q_j=p_j=0)_{j\neq i}$ 的那个平面内的一个解。这个解的\trFloquet{}乘子是什么？它是否\trNondegenerate？
\item 我们假设所有的 $\alpha_i$ 都相等。关于诸\trPeriodic{}解你能说些什么？
\end{enumerate}
\end{exercise}

\begin{exercise}\label{ex:ch14-010-lagrangian-j-orthogonal}
设 $E$ 是一个\trVectorSpace，它配有\trSymplecticForm{} $\omega$ 以及一个被 $\omega$ \trCalibrated{}的\trAlmostComplex{}结构 $J$；我们以 $\perp$ 表示对于\trInnerProduct{} $\omega(\cdot\,,J\cdot)$ 而言的\trOrthogonal{}性。证明：子空间 $L$ 是\trLagrange{}式的当且仅当 $L^\perp=JL$。
\end{exercise}

\begin{exercise}\label{ex:ch14-011-complex-submanifold}
设 $(W,\omega)$ 是一个\trSymplectic\trManifold，并设 $J$ 是一个被 $\omega$ \trCalibrated{}的\trAlmostComplex{}结构。设 $V\subset W$ 是一个复\trSubmanifold，也就是说，一个在 $J$ 之下稳定的\trSubmanifold：
\[
\forall\,x\in V\text{，}\quad J_x(T_xV)\subset T_xV\text{。}
\]
验证：$\omega$ 在 $V$ 上定义了一个\trNondegenerate{}的\trForm。
\end{exercise}

\begin{exercise}[复而不辛]\label{ex:ch14-012-hopf-surface}
考虑映射
\[
\begin{array}{rcl}
f\colon \CC^2 & \longrightarrow & \CC^2\\
(z_1,z_2) & \longmapsto & (2z_1,2z_2)
\end{array}
\]
以及 $\CC^2-\{0\}$ 对于由
\[
n\cdot(z_1,z_2)=f^n(z_1,z_2)
\]
给出的 $\ZZ$ 作用所成的商 $\scr{H}$（“\trHopf{}曲面”）。证明：这个商是一个复\trManifold（亦即其转移映射都是复解析的），并且\trDiffeomorphic{}于 $S^3\times S^1$。从而它的二阶\trRham{}\trCohomology{}群为零。由此推出：$\scr{H}$ 不具有任何辛结构。\footnote{在 [18]~中可以找到一些辛而不复的\trManifold{}的例子。}
\end{exercise}

\begin{exercise}[\trCayley{}数与\trAlmostComplex{}结构]\label{ex:ch14-013-cayley-almost-complex}
我们回忆：\trCayley{}数（Cayley numbers）或称\textbf{\trOctonion}（octaves）的代数是一个维数为 8 的实\trVectorSpace{} $\mathbf{O}$，它配有我们记作
\[
(1,i,j,k,\ell,\ell i,\ell j,\ell k)
\]
的一组基，并配有（对加法既左分配又右分配的）乘法，此乘法定义为
\[
i^2=j^2=k^2=\ell^2=ijk=jki=kij=-1\text{，}\qquad i\ell=-\ell i\text{，等等。}
\]

我们回忆：这个乘法既不交换也不结合（因为 $(ij)\ell=-i(j\ell)$），但它确实满足：对所有的 $a$ 与 $b$ 都有 $(ab)b=a(bb)$。

考虑由虚\trOctonion{}组成的\trEuclidean{}空间 $\RR^7$，也就是说，由 $(i,j,k,\ell,\ell i,\ell j,\ell k)$ 生成的那个子空间。设 $V\subset\RR^7$ 是一个维数为 6 的\trOrientation{}\trSubmanifold。对 $V$ 的任一点 $x$，我们以 $n(x)$ 表示单位\trNormalVector（由 $V$ 与 $\RR^7$ 的\trOrientation{}所定义），并如下定义
\[
J_x\colon T_xV\longrightarrow\RR^7\text{，}
\]
令 $J_x(u)=n(x)\cdot u$（\trOctonion{}意义下的乘法）。证明：$J_x$ 在 $T_xV$ 中取值，而且诸\trEndomorphism{} $J_x$ 在 $V$ 上定义了一个\trAlmostComplex{}结构 $J$。如此一来，$\RR^7$ 的所有\trHypersurface{}都是\trAlmostComplex{}的。\footnote{把 $\mathbf{O}$ 换成 $\mathbf{H}$、把 7 换成 3，便是用来证明“$\RR^3$ 中每个\trEmbedding{}的\trOrientation{}曲面都容有一个\trAlmostComplex{}结构”的一种类似（但更复杂）的途径。}
\end{exercise}

\begin{exercise}\label{ex:ch14-014-symplectic-blocks}
我们把有 $2n$ 行、$2n$ 列的矩阵写成分块矩阵
\[
A=\begin{pmatrix}X & Y\\ Z & T\end{pmatrix}\text{，}\qquad X,Y,Z,T\in\mathrm{M}_n(\RR)\text{。}
\]
在矩阵 $X$、$Y$、$Z$ 与 $T$ 满足哪些条件时，矩阵 $A$ 是群 $\mathrm{Sp}(2n;\RR)$ 的一个元素？
\end{exercise}

\begin{exercise}\label{ex:ch14-015-exp-symplectic}
设 $S$ 是一个实\trSymmetricMatrix。证明
\[
\exp(J_0S)\in\mathrm{Sp}(2n;\RR)\text{。}
\]
反之，假设矩阵 $S$ 使得
\[
\forall\,t\text{，}\quad\exp(tJ_0S)\in\mathrm{Sp}(2n;\RR)\text{；}
\]
关于 $S$ 你能说些什么？
\end{exercise}

\begin{exercise}\label{ex:ch14-016-determinant-one}
利用\cref{thm:sp-eigenspace-corollary}证明：\trSymplectic{}矩阵的行列式为 $1$。
\end{exercise}

\begin{exercise}[\trContact{}\trManifold{}的\trSymplectization]\label{ex:ch14-017-contact-symplectization}\BookIndexDefinition{T007}
设 $V$ 是一个\trManifold，它配有一个 1-\trForm{} $\alpha$，此形式是\trNonsingular{}的（也就是说，对每个 $x\in V$ 都有 $\alpha_x\neq 0$），并且满足
\[
\forall\,x\in V\text{，}\quad(d\alpha)_x|_{\operatorname{Ker}\alpha_x}\text{ 是一个\trNondegenerate{}的\trBilinearForm。}
\]
我们称 $\alpha$ 为 $V$ 上的一个\emph{\trContact{}形式}。$V$ 的维数必是奇数；我们把它写成 $2n+1$。

\begin{enumerate}[label=(\arabic*)]
\item 证明：$\alpha\wedge(d\alpha)^{\wedge n}$ 是 $V$ 上的一个\trVolume{}形式。
\item 证明：在 $V$ 上存在唯一一个\trVectorField{} $X$，使得
\[
i_X\alpha\equiv 1\qquad\text{且}\qquad i_X(d\alpha)=0
\]
（我们称 $X$ 为 $\alpha$ 的 \emph{\trReeb{}\trVectorField}）。
\item 考虑 $W=V\times\RR$，并给它配备由
\[
\omega_{(x,\sigma)}=d(e^{\sigma}\alpha)
\]
定义的 2-\trForm{} $\omega$。证明：$(W,\omega)$ 是一个\trSymplectic\trManifold（称为\trContact\trManifold{} $(V,\alpha)$ 的\emph{\trSymplectization}）。
\item 求出函数 $H(x,\sigma)=\sigma$ 的\trHamiltonian\trVectorField。
\end{enumerate}

\emph{例。}

设 $V=S^{2n+1}\subset\CC^{n+1}$ 为单位\trSphere（$\CC^{n+1}$ 配有坐标 $(q_1+ip_1,\dotsc,q_{n+1}+ip_{n+1})$），并考虑 1-\trForm
\[
\alpha=\frac{1}{2}\sum(p_i\,dq_i-q_i\,dp_i)\text{。}
\]

证明：$\alpha$ 是一个\trContact{}形式。求出 $\operatorname{Ker}\alpha$ 与 $\alpha$ 的\trReeb{}\trVectorField{} $X$。证明：$(S^{2n+1},\alpha)$ 的\trSymplectization{}与配有其标准\trSymplecticForm{}的 $\CC^{n+1}-\{0\}$ \trSymplectomorphism。
\end{exercise}

\begin{exercise}\label{ex:ch14-018-periodic-hamiltonian}
设 $H$ 是\trSymplectic\trManifold{}上的一个函数，若存在一个圆作用
\[
S^1\times W\longrightarrow W
\]
使得
\[
\frac{d}{dt}\Bigl(\exp(2i\pi t)\cdot x\Bigr)\Bigr|_{t=0}=X_H(x)
\]
对每个 $x\in W$ 都成立，我们就称 $H$ 为\trSymplectic\trManifold{}上的一个\emph{\trPeriodic\trHamiltonian{}量}\footnote{关于这个概念、一些基本结果以及更多内容，例如可参见 [5]。}。

\begin{enumerate}[label=(\arabic*)]
\item 证明：该圆作用的诸\trFixedPoint{}就是 $H$ 的诸\trCriticalPoint。
\item 证明：$H$ 的所有（\trPeriodic{}）\trTrajectory{}都是\trDegenerate{}的。
\end{enumerate}
\end{exercise}

\begin{exercise}\label{ex:ch14-019-circle-action-projective}
我们让\trCircle{} $S^1$ 经由
\[
u\cdot[z_0,\dotsc,z_n]=[u^{m_0}z_0,\dotsc,u^{m_n}z_n]\text{，}\quad m_i\in\ZZ
\]
作用在复\trProjectiveSpace{} $\CP^n(\CC)$ 上。证明：函数
\[
H([z_0,\dotsc,z_n])=\frac{1}{2}\,\frac{\sum m_i|z_i|^2}{\sum|z_i|^2}
\]
是与该作用相伴随的一个\trPeriodic\trHamiltonian{}量。\footnote{译注：本练习（\cref{ex:ch14-019-circle-action-projective}；连同\cref{ex:ch14-020-quadric-hamiltonians}、\cref{ex:ch14-021-isolated-critical-points}与\cref{ex:ch14-034-integral-periods}）照原书排印，惟归一化有可议之处。\cref{ex:ch14-018-periodic-hamiltonian}的定义以 $\exp(2i\pi t)$\trParametrization{}圆作用，即要求 $X_H$ 的时刻 1 流为恒等；而按第~5~章 $\CP^n(\CC)$ 上 $\sigma$ 的约定，本练习与\cref{ex:ch14-020-quadric-hamiltonians}所给\trHamiltonian{}量的\trHamiltonian{}流为 $t\mapsto[e^{im_i t}z_i]$（\cref{ex:ch14-021-isolated-critical-points}的局部正规形同理），时刻 1 并非恒等，须把\trHamiltonian{}量换成 $2\pi H$（或把圆作用改按 $e^{it}$\trParametrization）方能与该定义相合。又，\cref{ex:ch14-034-integral-periods}末段断言 $\int_{S^2}w^\star\omega\in\ZZ$，而按同一约定 $\CP^1(\CC)\subset\CP^n(\CC)$ 的面积大小为 $\pi$，故 $\int_{S^2}w^\star\omega$ 只能落在 $\pi\ZZ$ 之中，欲得 $\ZZ$ 须将 $\sigma$ 换成 $\sigma/\pi$。上述各处均照原书排印，存疑备案；归一化的取舍不影响各练习的定性结论。——译者注}

在（诸 $m_i$ 的）什么条件下，这个\trHamiltonian{}量是一个\trMorseFunction？此时它的诸\trCriticalPoint{}的\trIndex{}是什么？
\end{exercise}

\begin{exercise}\label{ex:ch14-020-quadric-hamiltonians}
我们回到第~1~章的\cref{ex:ch01-08}的那个\trQuadric{} $Q$。考虑函数 $g$ 与 $h$，即 $\CP^3(\CC)$ 上由
\[
g([z])=\frac{\operatorname{Im}(z_1\overline{z}_0)}{\sum|z_i|^2}
\qquad\text{与}\qquad
h([z])=\frac{\operatorname{Im}(z_3\overline{z}_2)}{\sum|z_i|^2}
\]
定义的函数在 $Q$ 上的限制。证明：$g$ 与 $h$ 是 $Q$ 上的\trPeriodic\trHamiltonian{}量。

如第~1~章的\cref{ex:ch01-08}中那样，我们取定实数 $\lambda$ 与 $\mu$，使得 $0<\lambda<\mu$，并考虑函数 $f$，即
\[
f([z])=\frac{\lambda\operatorname{Im}(z_1\overline{z}_0)+\mu\operatorname{Im}(z_3\overline{z}_2)}
             {\sum|z_i|^2}
\]
在 $Q$ 上的限制。

用 $g$ 与 $h$ 各自的\trHamiltonian\trVectorField{}表示出 $f$ 的\trHamiltonian\trVectorField{} $X_f$。由此推出 $f$ 的诸\trCriticalPoint。它们的\trIndex{}是什么（利用\cref{ex:ch14-018-periodic-hamiltonian}与\cref{ex:ch14-021-isolated-critical-points}，这个问题可以不作任何计算而得到回答）？
\end{exercise}

\begin{exercise}[难]\label{ex:ch14-021-isolated-critical-points}
考虑\trCompact{}\trSymplectic\trManifold{} $W$ 上的一个\trPeriodic\trHamiltonian{}量 $H$。设 $x$ 是 $H$ 的一个\trCriticalPoint。

\begin{enumerate}[label=(\arabic*)]
\item 利用\trDarboux{}定理的一个等变版本，证明：在 $x$ 的一个邻域内存在一个\trAlmostComplex{}结构与若干局部坐标，使得：
\begin{itemize}
\item $S^1$ 作用在复\trVectorSpace{} $T_xW$ 上的\trLinearization{}形如
\[
t\cdot(z_1,\dotsc,z_n)=\bigl(t^{\alpha_1}z_1,\dotsc,t^{\alpha_n}z_n\bigr)\text{。}
\]
\item 该\trHamiltonian{}量可以写成
\[
H(x_1,\dotsc,x_n;y_1,\dotsc,y_n)=\frac{1}{2}\sum_i\alpha_i\bigl(x_i^2+y_i^2\bigr)\text{。}
\]
\end{itemize}
\item 假设\trCriticalPoint{} $x$ 是\trIsolated{}的。证明：它是\trNondegenerate{}的并且\trIndex{}为偶数。
\item 证明：配备了\trPeriodic\trHamiltonian{}量、并且其诸\trCriticalPoint{}都是\trIsolated{}的\trConnected\trSymplectic\trManifold{}是\trSimplyConnected{}的。
\item 证明：\trCompact\trConnected\trSymplectic\trManifold{}上一个诸\trCriticalPoint{}全都\trIsolated{}的\trPeriodic\trHamiltonian{}量必有一个局部极小值与一个局部极大值。
\end{enumerate}
\end{exercise}

